\section*{Chapter \ref{ch:b3:schrodinger-three-dimensions} --- The Schrödinger Equation in Three Dimensions}

\begin{solution}{exo:b3:schrodinger-three-dimensions:1}
(a) $\int|\psi|^2\dd^3r = |A|^2\,4\pi\int_0^\infty r^2\eu^{-r^2/2
\sigma^2}\dd r = |A|^2(2\pi\sigma^2)^{3/2}$: $A = (2\pi\sigma^2)^{
-3/4}$. (b) By isotropy $\langle r^2\rangle = 3\langle x^2\rangle =
3\sigma^2$, so $\Delta x = \sigma$. (c) $\vect\jmath = \vect 0$: the
wave function is real. (d) $\langle\vect p\rangle = \hbar\vect k_0$
and $\vect\jmath = \rho\,\hbar\vect k_0/m$: the same cloud, now
drifting.
\end{solution}

\begin{solution}{exo:b3:schrodinger-three-dimensions:2}
$3\,(1)$, $6\,(3)$, $9\,(3)$, $11\,(3)$, $12\,(1)$, $14\,(6)$ ---
then nothing until $17$. Level $27$: $(3,3,3)$ and the permutations
of $(5,1,1)$ --- four states, since $27 = 9+9+9 = 25+1+1$. Cubic
symmetry only permutes axes, and no permutation links $(3,3,3)$ to
$(5,1,1)$: the coincidence is arithmetic (two representations as a
sum of three squares), not geometric.
\end{solution}

\begin{solution}{exo:b3:schrodinger-three-dimensions:3}
(a) $\vect\jmath = |A|^2\hbar k/m\,\vect e_x$. (b) Zero: real
function. (c) $j = (|A|^2\hbar k/m)(1 - r^2)$: incident flux minus
reflected flux --- the cross terms cancel. (d) $j = \rho\,\hbar
m/Mb$ with $\rho = 1/2\pi$ (per unit angle): a permanent
circulation, the microscopic ancestor of persistent currents.
\end{solution}

\begin{solution}{exo:b3:schrodinger-three-dimensions:4}
The coefficient is $\pi^2(\hbar c)^2/2\mu c^2 = \qty{3.76}{eV.nm^2}$.
(a) At $R = \qty{3}{nm}$: \qty{0.42}{eV}. (b) $E = \qty{2.16}{eV}$:
$\lambda = 1240/2.16 = \qty{575}{nm}$, yellow-green. (c) Smaller $R$,
larger $E$: toward the blue. (d) In bulk the confinement term
vanishes: fixed gap \qty{1.74}{eV}, $\lambda = \qty{713}{nm}$, the
same deep red whatever the sample.
\end{solution}

\begin{solution}{exo:b3:schrodinger-three-dimensions:5}
(a) Single-valuedness: $\psi(\varphi + 2\pi) = \psi(\varphi)$ forces
$\eu^{2\pi\iu m} = 1$. (b) $E_m = \hbar^2m^2/2Mb^2$; the states $\pm
m$ circulate oppositely and are exchanged by mirror reflection (or
time reversal) --- a symmetry of the ring, hence the degeneracy. (c)
The current loop $I = q\hbar m/2\pi Mb^2$ carries the magnetic moment
$\mu = I\pi b^2 = m\,q\hbar/2M$ --- quantised in steps of $q\hbar/2M$,
the Bohr magneton pattern of \cref{ch:b3:quantum-angular-momentum}.
(d) $\hbar^2/2m_{\text{e}}b^2 = \qty{1.9}{eV}$; the six electrons
fill $m = 0, \pm1$, and the first excitation $\pm1 \to \pm2$ costs
$(4-1) \times 1.9 \approx \qty{5.8}{eV}$, i.e.\ $\lambda \approx
\qty{210}{nm}$ --- the right ultraviolet neighbourhood from a ring
and nothing else.
\end{solution}

\begin{solution}{exo:b3:schrodinger-three-dimensions:6}
(a) Choose $n_x = 0..n$ and $n_y = 0..n - n_x$: $\sum_{n_x}(n - n_x
+ 1) = (n+1)(n+2)/2$. (b) With spin: $2, 6, 12, 20$; cumulative $2,
8, 20, 40$. (c) $2$, $8$, $20$ are magic; the failure beyond says
the nuclear well is not harmonic --- flatter-bottomed, and above all
possessed of a strong spin--orbit coupling that reshuffles the
shells to $28, 50, 82$. (d) Shell structure needs \emph{bunched}
levels separated by gaps; the box's levels spread irregularly with
no large gaps, so no nucleus-like stability pattern emerges.
\end{solution}

\begin{solution}{exo:b3:schrodinger-three-dimensions:7}
(a) $u_n = \sqrt{2/R}\sin(n\pi r/R)$, $E_n = n^2\pi^2\hbar^2/2MR^2$.
(b) $E_1 = \pi^2(\hbar c)^2/2Mc^2R^2 = 384210/(2 \times 939 \times
25) = \qty{8.2}{MeV}$. (c) The right scale: nucleons in nuclei are
MeV-spaced quantum states, as their spectra show. (d) $|u_1|^2$
peaks at $r = R/2$: the radial probability is largest halfway out
(the $r^2$ of the volume element, hidden in $u = r\varphi$, pushes
the maximum off-centre), while the 1D box's $|\varphi|^2$ peaks at
the centre.
\end{solution}

\begin{solution}{exo:b3:schrodinger-three-dimensions:8}
(a) Set $N = 2N_{\text{states}}(E_{\text{F}})$ and invert the count.
(b) Copper: $E_{\text{F}} = \qty{7.1}{eV}$. (c) $v_{\text{F}} =
\sqrt{2E_{\text{F}}/m_{\text{e}}} = \qty{1.6e6}{m/s}$;
$T_{\text{F}} = E_{\text{F}}/k_{\text{B}} \approx \qty{8e4}{K}$. (d)
Electrons are identical fermions: the Pauli exclusion principle
(\cref{ch:b3:identical-particles}) admits at most two per orbital
state, so the crowd must stack up to electron-volt energies.
\end{solution}

\begin{solution}{exo:b3:schrodinger-three-dimensions:9}
(a) $k = \sqrt{2M(V_0 - |E|)}/\hbar$, $\kappa = \sqrt{2M|E|}/\hbar$;
continuity of $u'/u$ at $R$ gives $k\cot kR = -\kappa$. (b) As $|E|
\to 0$, $\kappa \to 0$: $\cot kR = 0$, $kR = \pi/2$, and then $V_0 =
\hbar^2k^2/2M = \pi^2\hbar^2/8MR^2$. (c) With $\mu c^2 =
\qty{469}{MeV}$, $R = \qty{2.1}{fm}$: $V_{0,\min} = \qty{23}{MeV}$.
(d) The wall $u(0) = 0$ makes the 3D $s$ state the analogue of an
\emph{odd} 1D state, which must fit a quarter wave inside the well:
a shallow well cannot, whereas in 1D the nodeless even state always
binds.
\end{solution}

\begin{solution}{exo:b3:schrodinger-three-dimensions:10}
(a) $\kappa = \sqrt{2\mu c^2B}/\hbar c = \sqrt{2 \times 469 \times
2.22}/197.3 = \qty{0.231}{fm^{-1}}$: tail length $1/\kappa =
\qty{4.3}{fm}$, twice the range of the force. (b) $k =
\sqrt{2\mu(V_0 - B)}/\hbar$ and $\cot kR = -\kappa/k$. (c) Iterating:
$kR = 1.83$, $k = \qty{0.87}{fm^{-1}}$, $V_0 - B = (\hbar ck)^2/\mu
c^2\cdot\tfrac12 \approx \qty{31}{MeV}$: $V_0 \approx \qty{34}{MeV}$.
(d) Inside: $\int_0^R\sin^2kr\,\dd r \approx \qty{1.19}{fm}$;
outside: $\sin^2(kR)/2\kappa \approx \qty{2.0}{fm}$: about $63\%$ of
the probability lies \emph{beyond} the well --- the deuteron mostly
inhabits the region where its binding force has already given out, a
pure quantum halo.
\end{solution}

\begin{solution}{exo:b3:schrodinger-three-dimensions:11}
(a) $\dd\langle x\rangle/\dd t = \int x\,\partial_t\rho\,\dd^3r =
-\int x\operatorname{div}\vect\jmath\,\dd^3r = \int j_x\,\dd^3r =
\langle p_x\rangle/m$ (parts, boundary terms vanishing). (b)
Differentiate $\langle\vect p\rangle = -\iu\hbar\int\psi^*\vect\nabla
\psi$, use the equation twice; the kinetic terms cancel by parts,
leaving $-\int|\psi|^2\vect\nabla V$. (c) $V$ at most quadratic:
then $\langle\vect\nabla V(\vect r)\rangle = \vect\nabla V(\langle
\vect r\rangle)$ exactly. (d) Behind a double slit the packet is two
lobes; $\langle\vect r\rangle$ glides down the symmetry axis, where
the particle essentially never lands: the average describes the
ensemble's centroid, not anybody's trajectory.
\end{solution}

\begin{solution}{exo:b3:schrodinger-three-dimensions:12}
(a) Kinetic: gradients of scale $1/a$ give $\hbar^2/2m_{\text{e}}
a^2$; potential: the cloud sits at distances $\sim a$, giving
$-ke^2/a$ (the exact coefficients happen to be $1$). (b) $\dd E/\dd
a = 0$ at $a = \hbar^2/m_{\text{e}}ke^2 = a_0$; $E(a_0) =
-ke^2/2a_0 = \qty{-13.6}{eV}$. (c) Below $a_0$ the kinetic cost
grows as $1/a^2$, faster than the $-1/a$ gain: localisation is
taxed by the uncertainty principle, and the atom floats at the
break-even size. (d) Replace $ke^2$ by $Gm_{\text{e}}m_{\text{p}} =
\qty{1.0e-67}{J.m}$: $a_{\text{grav}} = \hbar^2/m_{\text{e}}G
m_{\text{e}}m_{\text{p}} \approx \qty{1.2e29}{m}$ --- larger than
the observable universe: gravity is too feeble to close a quantum
orbit, and builds stars instead of atoms.
\end{solution}

\begin{solution}{pb:b3:schrodinger-three-dimensions:1}
\textbf{1.} $s$-wave: $u = rR$ obeys the free 1D equation inside the
sphere with $u(0) = u(R) = 0$: $u \propto \sin(\pi r/R)$, energy
$\pi^2\hbar^2/2\mu R^2$.
\textbf{2.} The ground energy of a box is positive and grows as the
box shrinks --- localisation always costs kinetic energy
(uncertainty principle); it can only add to the gap.
\textbf{3.} With $\pi^2(\hbar c)^2/2\mu c^2 = \qty{3.76}{eV.nm^2}$:
$0.04$, $0.24$, $0.94$, \qty{3.8}{eV} --- at $R \approx \qty{2}{nm}$
the correction rivals the gap itself.
\textbf{4.} Attraction lowers the pair's energy: emission shifts
red. It scales as $1/R$ against the confinement's $1/R^2$: for small
dots the box term wins and the neglect improves.
\textbf{5.} $1240/1.74 = \qty{713}{nm}$: at the red edge of vision.
\textbf{6.} One colour from $10^{15}$ emitters certifies that the
synthesis made them all the \emph{same size} --- monodispersity, the
chemical feat behind the physics.
\textbf{7.} $1.97$, $2.34$, \qty{2.70}{eV}.
\textbf{8.} Confinement $0.23$, $0.60$, \qty{0.96}{eV}: $R =
\sqrt{3.76/\Delta E}$ gives $4.0$, $2.5$, \qty{2.0}{nm}.
\textbf{9.} Diameters $8.1$, $5.0$, \qty{4.0}{nm}: roughly $13$,
$8$ and $7$ lattice spacings across.
\textbf{10.} The blue dot is smallest, where $E$ depends most
steeply on $R$: a one-lattice-plane error shifts the colour most ---
small dots are the least forgiving.
\textbf{11.} Differentiate: $\dd E/\dd R = -2\Delta E_{\text{conf}}/R
= -\pi^2\hbar^2/\mu R^3$; at $R = \qty{2.5}{nm}$: $2 \times 0.60/2.5
= \qty{0.48}{eV/nm} = \qty{480}{meV/nm}$.
\textbf{12.} $\pm 5\% = \pm\qty{0.125}{nm}$: $\Delta E = \pm
\qty{60}{meV}$, i.e.\ $\Delta\lambda = \lambda^2\Delta E/hc \approx
\pm\qty{14}{nm}$ --- a $\sim\qty{28}{nm}$ spread, right at the
display's tolerance: five per cent is the \emph{boundary} of
acceptable chemistry.
\textbf{13.} Because colour purity demands size control at the
single-lattice-plane level across $10^{15}$ particles --- the
controlled growth achieved by Ekimov, Brus and Bawendi, and cited by
the 2023 Nobel committee.
\textbf{14.} A dot can only emit at its own (lowest) gap; absorbing
requires at least that energy, so the pump must be bluer. The excess
relaxes as lattice vibrations: heat.
\textbf{15.} $1 - 460/630 = 27\%$ of each pump photon's energy heats
the screen.
\textbf{16.} Above the emitting level the dot's spectrum is dense
(many box states): absorption succeeds over a broad band; the
excitation then tumbles to the lowest excited state, and emission
happens from that single level: narrow.
\textbf{17.} All dots absorb the same ultraviolet happily
(broadband), while each size emits its own colour: one lamp, many
labels --- the inverse of dye chemistry.
\textbf{18.} $\tfrac43\pi(4.0)^3/(0.3)^3 \approx \num{e4}$ atoms:
ten thousand atoms sharing one set of discrete, hydrogen-like
levels --- ``artificial atom'' is earned.
\textbf{19.} The crystal boundary: the confining wall of the
semiconductor grain replaces the nucleus's attraction as the agent
that holds and quantises the electron.
\textbf{20.} $\pi^2(\hbar c)^2/2Mc^2R^2 = 384210/(2 \times 939
\times 9) \approx \qty{23}{MeV}$: nuclear physics speaks MeV because
femtometre boxes do.
\textbf{21.} For an electron the box estimate turns relativistic;
honestly, $E \sim \pi\hbar c/R \approx \qty{200}{MeV}$ --- yet beta
electrons emerge with a few MeV: electrons cannot be constituents of
nuclei, the argument that buried the old nuclear-electron model and
prepared the neutrino's invention.
\textbf{22.} $\pi^2\hbar^2/2ML^2 \approx \qty{2e-70}{J}$:
thirty-nine orders below thermal noise --- people are not quantum
confined.
\textbf{23.} $\pi^2\hbar^2/2m_{\text{e}}R^2 = m_{\text{e}}c^2$ at
$R = \pi\hbar/\sqrt2\,m_{\text{e}}c \approx \qty{0.9}{pm}$, the
Compton scale: there the confinement energy suffices to create
electron--positron pairs and the one-particle equation abdicates to
quantum field theory.
\textbf{24.} $E_{\text{conf}} \propto 1/MR^2$: lighter particles and
smaller boxes are more quantum, in one formula.
\textbf{25.} $E_{\text{g}} + \pi^2\hbar^2/2\mu R^2$ maps $4.0 \to$
red, $2.5 \to$ green, $2.0\,\unit{nm} \to$ blue: the
particle-in-a-box, tuned by chemists to the nanometre, lights the
shop window --- confinement quantisation as consumer electronics.
\end{solution}
